%==============================================================================
% Section 1: Introduction
%==============================================================================
\section{Introduction}
\label{sec:introduction}

Driven by the global push for sustainable transport and
decarbonisation~\cite{iea2023}, lithium-ion batteries (LIBs) have become the
preferred energy-storage technology for electric vehicles (EVs) and grid-scale
renewable integration. Road transport causes about 15.9\,\% of global
greenhouse-gas emissions, mostly from passenger cars, and is a prime target for
electrification~\cite{crippa2024}; EU initiatives such as ``Fit for 55'' and
the ``Electric 2035'' vision further drive LIB demand, projected to grow almost
sevenfold over 2022--2030~\cite{statista2026}. However, complex ageing
mechanisms---solid-electrolyte interphase growth, active-lithium depletion and
electrode-structure deterioration---cause capacity fade and internal-resistance
rise, threatening the long-term reliability and safety of LIB
systems~\cite{meng2026,richardson2017}.

Battery health monitoring has shifted from physics-based methods (equivalent
circuit models, electrochemical simulation) to data-driven deep
learning~\cite{ghafari2023}. Recurrent architectures such as LSTM model
degradation trajectories well but suffer from a recurrence bottleneck on long
sequences~\cite{yao2021,zhao2024}. Transformer self-attention addresses this
but is quadratic, $O(L^{2})$, limiting real-time
deployment~\cite{patel2025}. Mamba and selective state space models break this
trade-off, matching Transformer accuracy at $O(L)$ linear
complexity~\cite{gu2023}. A key gap remains, however: most Mamba-based battery
studies are unimodal, using only macro-telemetry and overlooking
frequency-domain spectral features that are sensitive to the internal
degradation fingerprint~\cite{alsmadi2025,meng2026,olalde2024}.

Moreover, existing methods mostly estimate SOH and detect anomalies
separately, which is insufficient for real-time detection of anomalies from
transient events or sensor faults~\cite{arbaoui2024}, and few effective
conditioning mechanisms bridge the distribution gap between high-dimensional
temporal sequences and low-dimensional spectral features~\cite{meng2026}.
Building on this, we propose MambaSOHPredictor, a selective state space model
conditioned on physics-informed spectral features through Feature-wise Linear
Modulation (FiLM). Amenable to edge deployment, it supports efficient, robust
SOH prediction with simultaneous Isolation Forest anomaly detection for
reliable real-time battery health management.
